\documentclass[11pt,a4paper]{article}
\usepackage[margin=25mm]{geometry}
\usepackage{amsmath,amssymb,bm,microtype,newtxtext,newtxmath}
\usepackage[hidelinks]{hyperref}
\newcommand{\T}{^{\mathsf T}}
\title{Single-Sided and Fused Selected-Space CI Transport Contractions}
\author{Theory and implementation note}
\date{16 October 2026}
\begin{document}
\maketitle

\section{Selected space and required derivatives}

Let $W$ denote the fixed expansion of the selected VB structures into unique
spin-string products. The orthogonal-carrier image of a selected coefficient
vector $c$ is
\begin{equation}
 F(R)c=\operatorname{vec}\left[T_\alpha(R)X_c T_\beta(R)\T\right],
 \qquad \operatorname{vec}(X_c)=Wc,
\end{equation}
where $T_\sigma(R)=\wedge^{n_\sigma}R$, and $R$ is the upper-triangular active
orbital factor. The matrix $X_c$ is a mathematical representation only: the
implementation reads its touched columns through the sparse structure map.
Vectorization uses the fixed product index $aU_\beta+b$, where $a$ and $b$
label alpha and beta strings, respectively.
No selected-space closure under $T_\alpha\otimes T_\beta$ is assumed.

For two independent orbital factors, define
\begin{equation}
 f(L,R;\lambda,c)=\lambda\T F(L)\T F(R)c.
\end{equation}
The left and right cotangents are independent derivatives of this scalar.
A caller requesting only one must not compute the other. The identity
\begin{equation}
 f(L,R;\lambda,c)=f(R,L;c,\lambda)
\end{equation}
allows a left-only contraction to use the right-only kernel with both charts,
coefficients and supplied chart directions exchanged. It holds for any real
selected structure map $W$, including a nonclosed subspace.

At an accepted point $L=R=R_0$, the transport curvature consumes three objects:
\begin{align}
 G_d&=D_R\nabla_R f(R_0,R_0;\lambda,c)[\dot R],\\
 G_q&=\nabla_R f(R_0,R_0;\lambda,q),\\
 t&=DF(R_0)[\dot R]\T F(R_0)\lambda.
\end{align}
Here $\lambda$ is the accepted residual inverse image, $c$ the accepted CI
vector, $q$ its independent direction, and $\dot R$ the orbital direction.
The left chart is held fixed in $G_d$. None of these terms is discarded or
approximated. The mixed-block reciprocity is
\begin{equation}\label{eq:duality}
 \langle G_q,\dot R\rangle=q\T t.
\end{equation}

\section{Common tile intermediates}

All exterior matrices and directions below are evaluated at $R_0$. For a
tile $J$ of original beta strings, let $E_J$ select their unit vectors. Set
\begin{align}
 U&=T_\alpha X_cE_J,& V&=T_\beta E_J,& U_q&=T_\alpha X_qE_J,\\
 A&=T_\alpha\T U,& B&=T_\beta\T V,& A_q&=T_\alpha\T U_q,\\
 Z_\alpha&=T_\alpha X_\lambda B,&
 Z_\beta&=T_\beta X_\lambda\T A,&
 Z_{\beta q}&=T_\beta X_\lambda\T A_q.
\end{align}
Products with $X_\lambda$ use only its original touched products. The same
$V$, $B$ and $Z_\alpha$ serve the accepted and independent-CI directions.
Since only the right chart moves,
\begin{align}
 \dot U&=\dot T_\alpha X_cE_J,& \dot V&=\dot T_\beta E_J,\\
 \dot A&=T_\alpha\T\dot U,& \dot B&=T_\beta\T\dot V,\\
 \dot Z_\alpha&=T_\alpha X_\lambda\dot B,&
 \dot Z_\beta&=T_\beta X_\lambda\T\dot A.
\end{align}
The fixed left exterior factors are not differentiated in these expressions.

Define $\mathcal P_\sigma(Z,Y)$ as the upper-chart cotangent of
$\langle Z,T_\sigma C\rangle$ with respect to its right factor, where
$Y=T_\sigma C$. Its directional implementation differentiates the seed,
image and chart pullback, including the inverse-factor derivative already
present in the exterior generator pullback. The required contractions are
\begin{align}
 G_d&=\sum_J D\mathcal P_\alpha(Z_\alpha,U)
       [\dot Z_\alpha,\dot U,\dot R]
       +D\mathcal P_\beta(Z_\beta,V)[\dot Z_\beta,\dot V,\dot R],\\
 G_q&=\sum_J\mathcal P_\alpha(Z_\alpha,U_q)
                 +\mathcal P_\beta(Z_{\beta q},V).
\end{align}
Only these right cotangents are evaluated; their unused left counterparts
do not belong to either result.

\section{Eliminating the separate transpose-action sweep}

The raw coefficient-space tile of the transpose directional image is
\begin{equation}\label{eq:transpose}
 (t_{\mathrm{raw}})_J=\dot T_\alpha\T Z_\alpha
                         +T_\alpha\T\dot Z_\alpha.
\end{equation}
Both seeds on the right-hand side have already been obtained for $G_d$.
An exterior directional-adjoint action therefore produces this tile without
another expansion of $\lambda$ or an independent overlap-direction sweep.
For clarity, expanding the seeds gives
\begin{align}
 (t_{\mathrm{raw}})_J={}&\dot T_\alpha\T T_\alpha X_\lambda
                                  T_\beta\T T_\beta E_J\\
 &+T_\alpha\T T_\alpha X_\lambda T_\beta\T\dot T_\beta E_J.
\end{align}
These are precisely the alpha and beta contributions to
$DF(R_0)[\dot R]\T F(R_0)\lambda$. The final result is
\begin{equation}
 t=W\T\operatorname{vec}(t_{\mathrm{raw}}).
\end{equation}
Only entries touched by $W$ are sampled from each tile. This is an exact
projection, not a completion of the selected space. Equation~\eqref{eq:duality}
also follows immediately and provides an independent transpose check.

\section{Memory and unchanged numerical model}

At most eight alpha and seven beta row panels coexist. For width $b$ their
payload is bounded by $8b(8U_\alpha+7U_\beta)$ bytes. The planner uses the
conservative bound $64b(U_\alpha+U_\beta)$ within one 8 MiB panel allowance,
or the unavoidable one-column working set. One shared touched-product vector
contains $W\lambda$; $Wc$ and $Wq$ are generated only for the current beta
tile. The forcing is accumulated directly into the selected coefficients.
This factor order has no transformed spin-product image, pair-squared Gram,
or worker-owned replica of these arrays. Workers own disjoint row ranges of
shared panels.

When the touched support is the complete product carrier, the existing
one-body-moment order can be less expensive than the factor order. Its fused
implementation retains three shared product images only if they fit the
existing 64 MiB overlap allowance, not a larger workspace. These are
$L=F\lambda$, $Y=Fc$, and $\dot Y=DF[\dot R]c$. The right cotangent and its
direction use the one-body moments of $(L,Y)$ and $(L,\dot Y)$. After this
contraction, the storage of $Y$ is reused for $Fq$ and released. The dead
$\dot Y$ storage then holds the transpose directional image obtained by
applying the directional exterior adjoints to $L$. Thus the broad-support
order also shares its accepted image and avoids an additional overlap sweep.
If these three images do not fit, the same exact operator uses the bounded
factor sweep. This selects storage and contraction order, not the physical
state space; neither selected-space closure nor an approximation is invoked.

The full physical transport curvature, selected metric solves, orbital
gradient/HVP pullbacks and outer convergence criteria remain unchanged.
No residual threshold, rank cutoff, amplitude cutoff or CI-state approximation
is introduced. Tests compare the fused contraction with independently
evaluated cross-Gram derivatives, explicit occupied-minor finite differences,
and Eq.~\eqref{eq:duality}, including unequal spins, restricted supports,
nonclosed selections, zero directions and singular occupied minors.

\section{Finite-keyframe pullback in one sweep}

Let $R$ and $R_0$ denote the current and reference charts, respectively, and
let $\dot R_0$ be the accumulated reference direction. For a single state, write
\begin{equation}
 y_\lambda=F(R)\lambda,\quad y_c=F(R)c,\quad
 y_0=F(R_0)c_0,\quad
 z=F(R_0)b+DF(R_0)[\dot R_0]c_0-y_c.
\end{equation}
Here $b$ is the reference coefficient plus its independent CI displacement;
$c$ is the projected transported coefficient, not an assumption that the
target lies in the selected space. With $\lambda$ and these coefficients held
fixed, the two cotangents consumed by the finite transport are
\begin{align}
 \langle G_R,\delta R\rangle
 &=\langle DF(R)[\delta R]\lambda,z\rangle
   -\langle y_\lambda,DF(R)[\delta R]c\rangle,\\
 \langle G_0,\delta R_0\rangle
 &=\langle y_\lambda,DF(R_0)[\delta R_0]c_0\rangle.
\end{align}
The first line contains \emph{both} sides of the current metric derivative.
The second is the cotangent of the reference direction used in the transport
map; it does not differentiate the reference chart of the entire first line.
The CI-gradient consumer additionally requires
\begin{equation}
 t=F(R_0)\T y_\lambda.
\end{equation}

For an original beta tile $J$, form
\begin{align}
 U_0&=T_{0\alpha}X_{c_0}E_J,&V_0&=T_{0\beta}E_J,\\
 \dot U_0&=\dot T_{0\alpha}X_{c_0}E_J,&\dot V_0&=\dot T_{0\beta}E_J,\\
 U_b&=T_{0\alpha}X_bE_J+\dot U_0,&
 U_c&=T_\alpha X_cE_J,&V_c&=T_\beta E_J.
\end{align}
The target is the sum of $U_bV_0\T$, $U_0\dot V_0\T$, and
$-U_cV_c\T$. For any of these pairs $(U,V)$, common current-chart seeds are
\begin{equation}
 Z_\alpha=T_\alpha X_\lambda T_\beta\T V,\qquad
 Z_\beta=T_\beta X_\lambda\T T_\alpha\T U.
\end{equation}
The first pair's $Z_\alpha$ and the second pair's $Z_\beta$ also supply
$G_0$. Applying $T_{0\alpha}\T$ to that first $Z_\alpha$ and sampling only
the original touched entries supplies $t$. Thus the reference cotangent and
CI image need no independent cross-Gram sweep.

The factor planner conservatively allows sixteen alpha-plus-beta panels per
tile within the existing 64 MiB overlap allowance. There is one shared
$W\lambda$ vector, no pair-squared Gram or worker-owned panel replicas.
Complete broad support retains a moment contraction when five shared product
images fit that same allowance; otherwise the identical operator is streamed.

\section{Fixed-state energy and adjoint sharing}

For $m$ equally weighted prescribed states, define
\begin{equation}
 G=C\T M C,\qquad K=C\T H C,\qquad
 E=w\operatorname{Tr}(G^{-1}K),\qquad w=1/m.
\end{equation}
Only $MC$ is needed before normalization. Factor $G=L L\T$, set
$V=L^{-\mathsf T}$, and let $\widetilde C=CV$. A single shared sigma/adjoint
sweep produces $H\widetilde C$ and the unshifted integral cotangents. With
$\widetilde K=\widetilde C\T H\widetilde C$, the fixed-coefficient variation is
\begin{equation}
 \delta E=w\operatorname{Tr}(\widetilde C\T\delta H\widetilde C)
 -w\operatorname{Tr}(\widetilde K\widetilde C\T\delta M\widetilde C).
\end{equation}
The complete multiplier $-\widetilde K$, including its off-diagonal entries,
therefore completes the overlap cotangent after the same sweep. No second
Hamiltonian action is necessary. A small state-space Ritz transformation may
subsequently provide individual energies and representatives; equal-weight
energy and orbital cotangents are invariant under that transformation.
The derivative in the original coefficient frame is
\begin{equation}
 \nabla_C E=2w\big(H\widetilde C-M\widetilde C\widetilde K\big)V\T.
\end{equation}
All matrices inverted or diagonalized in this normalization are $m$ by $m$,
not full structure-space matrices. The implementation is in
\texttt{assembly/states.*}; bounded finite-transport contractions are in
\texttt{assembly/transport\_tiles.cpp}. Neither operation changes the selected
VB space, its spin carrier, the optimizer strategy, or its accuracy criteria.

\end{document}
